\documentclass[12pt]{article}

\usepackage{amsmath, amssymb}
\usepackage{graphicx}
\usepackage{geometry}
\usepackage{float}

\geometry{margin=1in}

\title{HW3}
\author{
Greyson Knippel \\[4pt]
MATH 339 --- Gonzaga University \\[4pt]
compiled with \LaTeX{}
}
\date{\today}

\begin{document}

\maketitle

\section*{Section 1.6 \#37}

answer the questions about the linear system
\[
\begin{cases}
-x \phantom{{}-3y} - z = -1 \\
2x \phantom{{}-3y} + 2z = 1 \\
x - 3y - 3z = 1
\end{cases}
\]

\subsection*{(a)}

\[
\begin{aligned}
-1x + 0y - 1z &= -1 \\
2x + 0y + 2z &= 1 \\
1x - 3y - 3z &= 1
\end{aligned}
\qquad\Longrightarrow\qquad
A = \begin{bmatrix} -1 & 0 & -1 \\ 2 & 0 & 2 \\ 1 & -3 & -3 \end{bmatrix}
\]

\subsection*{(b)}

finding the determinant:
\[
\det(A)
= -1\begin{vmatrix} 0 & 2 \\ -3 & -3 \end{vmatrix}
- 0\begin{vmatrix} 2 & 2 \\ 1 & -3 \end{vmatrix}
+ (-1)\begin{vmatrix} 2 & 0 \\ 1 & -3 \end{vmatrix}
\]
\[
= -1\bigl(0 - (-6)\bigr) - 0 + (-1)\bigl(-6 - 0\bigr)
= -6 + 6 = 0
\]

\subsection*{(c)}

nope. $\det(A) = 0$, so $A$ is not invertible, and $A\vec{x} = \vec{b}$ does not have a unique solution.

\subsection*{(d)}

row reduce the matrix
\[
\left[\begin{array}{rrr|r}
-1 & 0 & -1 & -1 \\
2 & 0 & 2 & 1 \\
1 & -3 & -3 & 1
\end{array}\right]
\xrightarrow{(-1)R_1}
\left[\begin{array}{rrr|r}
1 & 0 & 1 & 1 \\
2 & 0 & 2 & 1 \\
1 & -3 & -3 & 1
\end{array}\right]
\]
\[
\xrightarrow[R_3 - R_1]{R_2 - 2R_1}
\left[\begin{array}{rrr|r}
1 & 0 & 1 & 1 \\
0 & 0 & 0 & -1 \\
0 & -3 & -4 & 0
\end{array}\right]
\xrightarrow{R_2 \leftrightarrow R_3}
\left[\begin{array}{rrr|r}
1 & 0 & 1 & 1 \\
0 & -3 & -4 & 0 \\
0 & 0 & 0 & -1
\end{array}\right]
\]
row 3 says $0x + 0y + 0z = -1$, i.e.\ $0 = -1$, which is impossible. so it has no solutions.

\section*{Section 1.7 \#25}

find $A^{-1}$ using an $LU$ factorization, where
\[
A = \begin{bmatrix} 1 & 4 \\ -3 & -11 \end{bmatrix}
\]

row reduce $A$ to upper triangular form:
\[
\begin{bmatrix} 1 & 4 \\ -3 & -11 \end{bmatrix}
\xrightarrow{R_2 + 3R_1}
\begin{bmatrix} 1 & 4 \\ 0 & 1 \end{bmatrix} = U,
\qquad
L = \begin{bmatrix} 1 & 0 \\ -3 & 1 \end{bmatrix}
\]
we can check:
\[
LU = \begin{bmatrix} 1 & 0 \\ -3 & 1 \end{bmatrix}\begin{bmatrix} 1 & 4 \\ 0 & 1 \end{bmatrix}
= \begin{bmatrix} 1 & 4 \\ -3 & -11 \end{bmatrix} = A \quad \checkmark
\]

invert $L$ and $U$:
\[
L^{-1} = \frac{1}{(1)(1) - (0)(-3)}\begin{bmatrix} 1 & 0 \\ 3 & 1 \end{bmatrix}
= \begin{bmatrix} 1 & 0 \\ 3 & 1 \end{bmatrix},
\qquad
U^{-1} = \frac{1}{(1)(1) - (4)(0)}\begin{bmatrix} 1 & -4 \\ 0 & 1 \end{bmatrix}
= \begin{bmatrix} 1 & -4 \\ 0 & 1 \end{bmatrix}
\]

since $A = LU$, $A^{-1} = U^{-1}L^{-1}$:
\[
A^{-1} = \begin{bmatrix} 1 & -4 \\ 0 & 1 \end{bmatrix}\begin{bmatrix} 1 & 0 \\ 3 & 1 \end{bmatrix}
= \begin{bmatrix} (1)(1) + (-4)(3) & (1)(0) + (-4)(1) \\ (0)(1) + (1)(3) & (0)(0) + (1)(1) \end{bmatrix}
= \begin{bmatrix} -11 & -4 \\ 3 & 1 \end{bmatrix}
\]
and then check:
\[
AA^{-1} = \begin{bmatrix} 1 & 4 \\ -3 & -11 \end{bmatrix}\begin{bmatrix} -11 & -4 \\ 3 & 1 \end{bmatrix}
= \begin{bmatrix} -11 + 12 & -4 + 4 \\ 33 - 33 & 12 - 11 \end{bmatrix}
= \begin{bmatrix} 1 & 0 \\ 0 & 1 \end{bmatrix} = I \quad \checkmark
\]

\end{document}
